\subsection{Конфигурационный файл}

\subsubsection{Описание решения}

Сфокусировавшись на потоке ввода, упростим задачу до трёх ситуаций:
%
\begin{enumerate}
  \item Начало блока.
  \item Конец блока.
  \item Присваивание переменной значения.
\end{enumerate}

Введём «глобальную» хеш-таблицу, которая хранит текущее состояние переменных. Ключи --- имена переменных, значение по ключу --- значение переменной.

Также введём стек, элементы которого будут эквивалентны вложенным блокам. Идея --- хранить в нём историю изменений значений переменных. Вершина стека хранит значения переменных на прошлом уровне.

Когда из потока ввода поступает команда «начать новый блок», в стек будет записана пустая хеш-таблица. По мере обновления значений переменных в текущий блок будет записываться старое значение переменной, а в глобальную хеш-таблицу --- актуальное значение.

Как только поступит команда «завершить текущий блок», из вершины стека удаляется блок. Затем, для каждой переменной из удаляемого блока мы восстанавливаем старое значение в глобальной хеш-таблице. Так мы восстанавливаем исходное состояние переменных до входа в очередной блок.

\subsubsection{Асимптотическая сложность}

\textit{Временная сложность решения} --- линейная \( \mathcal{O}(N) \). Для каждой считанной строки выполняется константное количество операций в случае начала блока и обновления значения переменной. Если же блок завершается, то все предыдущие присвоения блока отменяются за константное время --- в сумме такие операции не превысят линейное время работы.

\textit{Пространственная сложность решения} --- линейная \( \mathcal{O}(N) \). Для хранения текущего состояния используется стек, элементами которого являются хеш-таблицы. Так как и добавление элемента в стек, и добавление элемента в хеш-таблицу происходит по строчной команде, число всех элементов в сумме не превышает числа входных строк.

\subsubsection{Листинг кода}

\begin{lstlisting}
#include <bits/stdc++.h>

#include <iostream>
#include <stack>
#include <unordered_map>

using scope = std::unordered_map<std::string, long>;

void update_scope(std::stack<scope>& ctx, scope& vars, std::string lval, std::string rval) {
  char* fail;
  long rval_long = std::strtol(rval.c_str(), &fail, 10);
  if (ctx.top().count(lval) == 0) {
    ctx.top()[lval] = vars.count(lval) ? vars[lval] : 0;
  }
  if (*fail) {
    vars[lval] = vars[rval];
    std::cout << vars[lval] << "\n";
  } else {
    vars[lval] = rval_long;
  }
}

int main() {
  std::string line;
  scope vars;
  std::stack<scope> ctx;
  ctx.push({});
  while (std::cin >> line) {
    if (line == "{") {
      ctx.push({});
    } else if (line == "}") {
      // insert_range does not suit there:
      // "if there is no element with key equivalent to the key of that element"
      for (const auto& [key, value] : ctx.top()) {
        vars[key] = value;
      }
      ctx.pop();
    } else {
      std::stringstream ss(line);
      std::string lval, rval;
      std::getline(ss, lval, '=');
      std::getline(ss, rval, '=');
      update_scope(ctx, vars, lval, rval);
    }
  }
  return 0;
}
\end{lstlisting}